\documentclass[../main.tex]{subfiles}
\begin{document}

\section{Reading and evidence}
\label{sec:kpc:interface}

\begin{definition}[Exported membership]
\label{def:kpc-interface:membership-reading}
For a handle $x$ the interface answers whether a record belongs to $C_k(x)$.
Handles are defined in \zcref{def:kpc-model:keyed-membership}; records and
their sets in \zcref{def:kpc-operations:classification}, and paths in
\zcref{def:kpc-operations:continuation}.
\end{definition}

\begin{corollary}[Reading along a path]
\label{cor:kpc-interface:path-reading}
For a handle $g$, at each prefix containing $g$ the live handle $y$ with
$g\preceq y$ is unique, and along extensions of one history the readings
$C_k(y)$ grow.
\end{corollary}
\begin{proof}
\zcref{thm:kpc-invariants:unique-successor} gives the unique live end at
each prefix. An earlier end $y$ lies on the same path as a later end $y'$,
so $y\preceq y'$ and $C_k(y)\subseteq C_k(y')$ by
\zcref{thm:kpc-invariants:set-membership-growth}; each set is fixed by its
dependencies (\zcref{thm:kpc-invariants:source-provenance}).
\end{proof}

\begin{remark}[Record evidence]
\label{rem:kpc-interface:record-evidence}
A record exposes its SOURCE handle's creating event; a reader's own
published recognition rule may read data of that event, such as what it
carries and to which party it is attributed, to decide which records it
accepts, while the covenant checks none of it.
Acceptance filters a reading and changes no set $C_k$; a fixed rule that
depends only on the SOURCE event's immutable contents keeps accepted
readings monotone along a path by \zcref{cor:kpc-interface:path-reading}.
\end{remark}

\paragraph{Named evidence assumptions.}
\emph{Inclusion} authenticates that the events and dependencies belong
to one valid history with the identity and single-consumption properties
of \zcref{sec:kpc:model}.
\emph{Retrieval} supplies the finite dependency closure, complete input
classification, each spent output's birth position and output order.
\emph{Finality} holds for a read handle once it lies beyond the depth the
realization publishes, a quantity of the realization and not of this
system: no history replacing the one that contains it is then accepted.
A reading claimed final requires it.
These assumptions concern access to authentic history and its permanence;
disjoint keyed membership and readable covenant clauses remain the
realization
assumptions of \zcref{def:kpc-model:keyed-membership}.

\paragraph{Reconstruction.}
Authenticate the key and the handle $x$. For each reached handle $s$,
authenticate its creating event $e$, first-output birth position, inclusion
and every input's membership. Check both covenant clauses for consumed
handles. If $I_k(e)=\varnothing$, return the singleton $\{r_k(s)\}$.
Otherwise recursively reconstruct all handles in $I_k(e)$ and return
their union. Successive readings may reuse completed handles.

Missing evidence yields \emph{incomplete}, not a set by default.
Conflicting identities or cyclic dependencies yield \emph{invalid}: they
establish no valid history. A later covenant-shaped output is not a handle
and has no set to read. These are evidence distinctions; the recurrence defines a set
for every handle of every valid history.

\begin{theorem}[Faithful reading]
\label{thm:kpc-interface:membership-correctness}
With inclusion and retrieval evidence, reconstruction terminates and
returns exactly $C_k(x)$ for every handle $x$.
\end{theorem}
\begin{proof}
Finite acyclic dependencies permit induction. At a leaf, complete input
classification establishes a SOURCE and the reader returns its defined
singleton. At a MERGE, classification supplies precisely $I_k(e)$ and
the recursively obtained input sets give exactly the defining union.
Complete evidence establishes every branch, so the return equals $C_k(x)$.
\end{proof}

\end{document}
